\documentclass[a4paper,12pt]{scrartcl}

\usepackage{color}
\usepackage{graphicx}
\usepackage{float}
\usepackage[utf8]{inputenc}

\setcounter{tocdepth}{2}

\usepackage{listings}

\definecolor{keyword}      {rgb}{0.7, 0.3, 0.2}
\definecolor{comment}      {rgb}{0.4, 0.7, 0.7}
\definecolor{string}       {rgb}{0.0, 0.7, 0.3}
\definecolor{line_numbers} {rgb}{0.8, 0.8, 0.8}
\lstset{
    showspaces=false, 
    showstringspaces=false,
    basicstyle=\ttfamily,
    keywordstyle=\color{keyword}, 
    commentstyle=\color{comment}, 
    stringstyle=\color{string},
    numbers=left, 
    numberstyle=\color{line_numbers}, 
    stepnumber=1, 
    numbersep=10pt,
    breaklines=true
}



\title{\LaTeX{}watchr}
\author{Patrick Oscity}
\date{\today}

\begin{document}

	\maketitle
	\tableofcontents
	\pagebreak

	\section{Introduction}
	Welcome to \LaTeX{}watchr. If you managed to compile this document,
	everything should be up and running\footnote{Congratulations!}. Now
	go ahead and start writing a fancy \LaTeX{} document. In the
	following sections you will find some content to test if everything
	works as expected.

	\section{UTF-8 Encoding}
	Just checking: ÄäÁáÀàËëÉéÈèÍíÌìÖöÓóÒòÜüÚúÙùß fi ffi fl ffl

	\section{Mathematical Equations}
	Now, let's have a look at the Quadratic Formula:
	\begin{equation}
		\begin{array}{*{20}c}
			{x = \frac{{ - b \pm \sqrt {b^2  - 4ac} }}{{2a}}} &
			{{\rm{when}}} & {ax^2  + bx + c = 0}  \\
		\end{array}
	\end{equation}

	\section{Source Code}
	\lstinputlisting[
	    language=Ruby,
	    label=Codebeispiel,
	    caption=Excerpt from the \LaTeX{}watchr source code
	]{../latexwatchr/latex.rb}
	
	\section{Just Some Text}
	Lorem ipsum dolor sit amet, consectetur adipiscing elit. Curabitur
	ut ante at risus malesuada lacinia consectetur laoreet justo. Aenean
	feugiat pellentesque velit, eget placerat justo porttitor at.
	Quisque adipiscing metus metus. Vivamus et vehicula nibh. Quisque
	consequat, augue in interdum accumsan, lorem leo bibendum elit, a
	accumsan augue dui eu est. Proin vel dolor ut neque commodo malesuada.
	\par
	Vestibulum ipsum diam, varius non euismod non, accumsan eget lectus.
	Nam ullamcorper, elit non facilisis ultricies, ipsum libero porttitor
	nisi, eu porttitor nisi arcu luctus dui. In quis nulla quam. Maecenas
	semper, arcu in tincidunt facilisis, felis diam ornare risus, sit amet
	mattis erat nisi eget nibh. In hac habitasse platea dictumst. Praesent
	lacinia mattis hendrerit. Donec sollicitudin lobortis vehicula.
	Integer imperdiet nulla eget enim vulputate et rhoncus massa
	venenatis. Praesent vehicula nulla id lorem euismod non interdum
	tortor vulputate. Quisque vel mauris mauris.

	\section{Embedding Images}
	Now, let's have a look at Figure ~\ref{myfig}. Is everything working here? I hope so…
	\begin{figure}[h]
	    \centering
	    \includegraphics[width=50mm]{images/sample_image.png}
	    \caption{sample image with color and transparency}
		\label{myfig}
	\end{figure}
\end{document}
